\section{Kiến thức nền tảng và công nghệ sử
dụng}\label{kiux1ebfn-thux1ee9c-nux1ec1n-tux1ea3ng-vuxe0-cuxf4ng-nghux1ec7-sux1eed-dux1ee5ng}

\subsection{NLP}\label{nlp}

\subsubsection{Bag of word}\label{bag-of-word}

Bag of Words (BoW) là một kỹ thuật biểu diễn văn bản đơn giản trong xử
lý ngôn ngữ tự nhiên, trong đó mỗi tài liệu được xem như một "túi" chứa
các từ mà không quan tâm đến thứ tự hay ngữ pháp. Mô hình này đếm số lần
xuất hiện của từng từ trong tập từ vựng để tạo thành vector đặc trưng,
giúp máy học có thể xử lý và so sánh các văn bản dựa trên tần suất xuất
hiện của từ. BoW dễ triển khai nhưng không nắm được ngữ cảnh hay ý nghĩa
của từ trong câu.

\includegraphics[width=4.96354in,height=2.14372in]{media/image1.png}

\emph{\textbf{Hình 3.1.1. Bag of words.}}

\emph{Nguồn: ``Bag of Words in NLP'', Medium. Trích từ :
https://ayselaydin.medium.com/4-bag-of-words-model-in-nlp-434cb38cdd1b}

\subsubsection{TF-IDF}\label{tf-idf}

TF-IDF (Term Frequency -- Inverse Document Frequency) là một phương pháp
biểu diễn văn bản nhằm đo lường mức độ quan trọng của một từ trong một
tài liệu so với toàn bộ tập tài liệu. TF thể hiện tần suất xuất hiện của
từ trong tài liệu, còn IDF giảm trọng số của những từ xuất hiện phổ biến
ở nhiều tài liệu vì chúng ít mang ý nghĩa phân biệt. Khi nhân TF và IDF,
ta thu được trọng số TF-IDF: giá trị càng cao cho thấy từ đó vừa xuất
hiện nhiều trong tài liệu, vừa ít gặp trong tài liệu khác, từ đó phản
ánh mức độ quan trọng của nó. TF-IDF thường được dùng trong tìm kiếm,
phân loại, và lọc thông tin văn bản.

\includegraphics[width=4.80279in,height=2.86412in]{media/image2.png}

\emph{\textbf{Hình 3.1.2. TF-IDF.}}

\emph{Nguồn: ``Text Vectorization and Word Embedding \textbar{} Guide to
Master NLP (Part 5)'', Analytics Vidhya. Trích từ :
https://www.analyticsvidhya.com/blog/2021/06/part-5-step-by-step-guide-to-master-nlp-text-vectorization-approaches/}

\subsubsection{Word Embeddings}\label{word-embeddings}

Word Embedding là một kỹ thuật biểu diễn từ trong không gian vector liên
tục, trong đó mỗi từ được ánh xạ thành một vector số thực có kích thước
cố định. Khác với Bag of Words hay TF-IDF chỉ dựa vào tần suất, Word
Embedding nắm bắt được ngữ nghĩa và mối quan hệ ngữ cảnh giữa các từ.
Nhờ đó, các từ có nghĩa gần nhau thường nằm gần nhau trong không gian
vector (ví dụ: \emph{king} gần \emph{queen}, \emph{cat} gần \emph{dog}).
Word Embedding là nền tảng quan trọng trong nhiều mô hình học sâu xử lý
ngôn ngữ tự nhiên như Word2Vec, GloVe, FastText và các mô hình
Transformer hiện đại.

\includegraphics[width=3.69758in,height=2.75565in]{media/image3.png}

\emph{\textbf{Hình 3.1.3. Word embeddings.}}

\emph{Nguồn: ``Word Embedding: Basics'', Medium. Trích từ :
https://medium.com/@hari4om/word-embedding-d816f643140}

\subsubsection{RNN}\label{rnn}

Recurrent Neural Network (RNN) là một loại mạng nơ-ron được thiết kế để
xử lý dữ liệu theo chuỗi, chẳng hạn như văn bản, âm thanh hoặc chuỗi
thời gian. Điểm đặc biệt của RNN là bộ nhớ ẩn (hidden state), cho phép
mạng ghi nhớ thông tin từ các bước trước và sử dụng nó để dự đoán bước
tiếp theo, nhờ đó nắm bắt được mối quan hệ tuần tự trong dữ liệu.

RNN thường được dùng trong các bài toán như dự đoán chuỗi, dịch máy,
nhận dạng giọng nói, và tạo văn bản. Tuy nhiên, RNN truyền thống gặp vấn
đề vanishing/exploding gradient khi học các chuỗi dài, dẫn đến khó khăn
trong việc nắm bắt ngữ cảnh xa. Các biến thể như LSTM (Long Short-Term
Memory) và GRU (Gated Recurrent Unit) đã được phát triển để giải quyết
vấn đề này, cho phép ghi nhớ thông tin lâu hơn và cải thiện hiệu suất
của mô hình.

\includegraphics[width=6.26772in,height=1.86111in]{media/image4.png}

\emph{\textbf{Hình 3.1.4. RNN.}}

\emph{Nguồn: ``What is Recurrent Neural Networks (RNN)?'', Analytics
Vidhya. Trích từ :
https://www.analyticsvidhya.com/blog/2022/03/a-brief-overview-of-recurrent-neural-networks-rnn/}

\subsubsection{Sequence-to-Sequence}\label{sequence-to-sequence}

Sequence-to-Sequence (Seq2Seq) là một kiến trúc mạng nơ-ron được thiết
kế để chuyển đổi một chuỗi đầu vào thành một chuỗi đầu ra, thường có độ
dài khác nhau. Kiến trúc này bao gồm encoder và decoder:

\begin{itemize}
\item
  Encoder đọc toàn bộ chuỗi đầu vào và mã hóa thông tin thành một vector
  ngữ cảnh (context vector).
\item
  Decoder sử dụng vector này để sinh ra chuỗi đầu ra từng bước một, dựa
  trên ngữ cảnh và các từ đã sinh trước đó.
\end{itemize}

Seq2Seq được ứng dụng rộng rãi trong các bài toán như dịch máy (machine
translation), tóm tắt văn bản (text summarization), trả lời tự động
(chatbots) và nhiều tác vụ chuyển đổi chuỗi khác. Khi kết hợp với cơ chế
attention, Seq2Seq có thể nắm bắt ngữ cảnh tốt hơn, cải thiện chất lượng
dự đoán cho các chuỗi dài.

\includegraphics[width=4.92708in,height=1.55208in]{media/image5.png}

\emph{\textbf{Hình 3.1.5. Seq2seq.}}

\emph{Nguồn: ``Sequence-to-Sequence Learning for Machine Translation'',
Dive Into Deep Learning. Trích từ :
https://d2l.ai/chapter\_recurrent-modern/seq2seq.html}

\subsubsection{Attention Mechanism}\label{attention-mechanism}

Attention mechanism trong RNN giúp mô hình tập trung vào những phần quan
trọng của chuỗi đầu vào khi tạo ra mỗi bước của đầu ra. Thay vì chỉ dựa
vào hidden state cuối cùng (vốn dễ mất thông tin khi chuỗi dài),
attention tạo ra một trọng số cho từng hidden state của encoder, cho
biết mức độ liên quan của nó đối với bước giải mã hiện tại. Sau đó, mô
hình tạo ra context vector bằng cách nhân các hidden state với trọng số
attention này. Nhờ đó, RNN có thể ``nhìn lại'' toàn bộ chuỗi đầu vào và
chọn lọc thông tin quan trọng, giúp cải thiện đáng kể chất lượng dịch,
tóm tắt văn bản và các tác vụ seq2seq khác.

\includegraphics[width=6.26772in,height=2.91667in]{media/image6.png}

\emph{\textbf{Hình 3.1.6. Attention Mechanism.}}

\emph{Nguồn: ``The Attention Mechanism Before Transformers'', The
AIedge. Trích từ :
https://newsletter.theaiedge.io/p/the-attention-mechanism-before-transformers}

\subsection{Nghiên cứu LLM}\label{nghiuxean-cux1ee9u-llm}

\subsubsection{Nghiên cứu cách máy tính hiểu ngôn ngữ tự
nhiên}\label{nghiuxean-cux1ee9u-cuxe1ch-muxe1y-tuxednh-hiux1ec3u-nguxf4n-ngux1eef-tux1ef1-nhiuxean}

\paragraph{Tokenization}\label{tokenization}

Tokenization là quá trình chuyển đổi văn bản tự nhiên thành các đơn vị
nhỏ hơn gọi là token. Một token có thể là một từ, một ký tự, hoặc một
đoạn văn bản phụ thuộc vào phương pháp tokenization được sử dụng.

Trong các mô hình ngôn ngữ lớn (LLM), token thường là các phần từ con
(subword) nhờ vào các kỹ thuật như Byte Pair Encoding (BPE) hoặc
SentencePiece.

Quá trình này giúp mô hình hiểu và xử lý ngôn ngữ hiệu quả hơn, đồng
thời giảm số lượng từ vựng cần thiết để huấn luyện, đặc biệt hữu ích khi
làm việc với nhiều ngôn ngữ hoặc từ ngữ hiếm.

\includegraphics[width=3.97344in,height=2.23438in]{media/image7.png}

\emph{\textbf{Hình 3.2.2.1. Chuyển một câu thành các token nhỏ hơn.}}

\emph{Nguồn: ``Tokenization --- A complete guide'', Medium. Trích từ :
https://medium.com/@utkarsh.kant/tokenization-a-complete-guide-3f2dd56c0682}

\paragraph{Embedding}\label{embedding}

Embedding là quá trình chuyển các token sau khi tokenization thành các
vector số trong không gian nhiều chiều, nhằm giúp mô hình xử lý được
ngôn ngữ tự nhiên dưới dạng dữ liệu số.

Mỗi token sẽ được ánh xạ thành một vector có giá trị thực, thể hiện ngữ
nghĩa và mối quan hệ giữa các từ trong ngữ cảnh. Các phương pháp phổ
biến như Word2Vec, GloVe, hay embedding học được từ các mô hình
transformer như BERT hoặc GPT giúp mô hình hiểu được sự tương đồng ngữ
nghĩa và ngữ cảnh của từ ngữ. Embedding là bước quan trọng để máy tính
``hiểu'' được ngôn ngữ con người.

\includegraphics[width=5.64678in,height=3.07665in]{media/image8.png}

\emph{\textbf{Hình 3.2.1.2. Chuyển một câu thành một vector.}}

\emph{Nguồn: A Complete Guide to Embedding For NLP \& Generative AI/LLM,
Toward AI. Trích từ:
https://pub.towardsai.net/a-complete-guide-to-embedding-for-nlp-generative-ai-llm-4a24301aba97}

\subsubsection{Nghiên cứu quy trình huấn luyện mô hình ngôn ngữ lớn
(LLM) từ
đầu}\label{nghiuxean-cux1ee9u-quy-truxecnh-huux1ea5n-luyux1ec7n-muxf4-huxecnh-nguxf4n-ngux1eef-lux1edbn-llm-tux1eeb-ux111ux1ea7u}

Huấn luyện một mô hình ngôn ngữ lớn (LLM) từ đầu là một quá trình phức
tạp và tốn nhiều tài nguyên, bao gồm ba bước chính: chuẩn bị dữ liệu,
chuẩn bị phần cứng, và thực hiện huấn luyện mô hình.

\paragraph{Chuẩn bị dữ liệu}\label{chuux1ea9n-bux1ecb-dux1eef-liux1ec7u}

Việc huấn luyện LLM đòi hỏi một lượng lớn dữ liệu văn bản đa dạng và
chất lượng cao. Dữ liệu có thể được thu thập từ các nguồn như sách,
trang web, tài liệu kỹ thuật,... Sau đó, dữ liệu cần được làm sạch, loại
bỏ các nội dung không phù hợp và chuẩn hóa định dạng để phù hợp với quá
trình tokenization.

\includegraphics[width=4.08233in,height=3.06175in]{media/image9.png}

\emph{\textbf{Hình 3.2.2.1. Sử dụng dữ liệu từ internet để huấn luyện
LLM.}}

\emph{Nguồn: Làm thế nào để đào tạo Large Language Models, VinBigData.
Trích từ:
https://vinbigdata.com/cong-nghe-giong-noi/lam-the-nao-de-dao-tao-large-language-models.html}

\paragraph{Chuẩn bị hardware}\label{chuux1ea9n-bux1ecb-hardware}

Huấn luyện LLM yêu cầu phần cứng mạnh, đặc biệt là các GPU hiệu năng cao
hoặc cụm máy tính phân tán. Ngoài ra, cần cân nhắc bộ nhớ GPU, băng
thông I/O và khả năng lưu trữ lớn để xử lý hàng tỷ token trong quá trình
huấn luyện.

\includegraphics[width=4.35573in,height=2.45282in]{media/image10.png}

\emph{\textbf{Hình 3.2.2.2. NVIDIA GPU 4090\\
}Nguồn: GeForce RTX 4090, Nvidia. Trích từ:
https://www.nvidia.com/en-gb/geforce/graphics-cards/40-series/rtx-4090/}

\paragraph{Training LLM}\label{training-llm}

Quá trình huấn luyện bao gồm việc truyền dữ liệu qua mạng neural
transformer, điều chỉnh trọng số qua các bước lan truyền ngược
(backpropagation) nhằm tối ưu hóa hàm mất mát. Thời gian huấn luyện có
thể kéo dài từ vài ngày đến vài tuần tùy vào kích thước mô hình và tài
nguyên phần cứng. Trong giai đoạn này, việc giám sát quá trình huấn
luyện và điều chỉnh siêu tham số là rất quan trọng để đạt được hiệu quả
tối ưu.

\includegraphics[width=4.81128in,height=2.62143in]{media/image11.png}

\emph{\textbf{Hình 3.2.2.3. Mô hình Transformer.}}

\emph{Nguồn: Foundation Models, Transformers, BERT and GPT, Niklas
Heidloff. Trích từ:
https://heidloff.net/article/foundation-models-transformers-bert-and-gpt/}

\subsubsection{Nghiên cứu quy trình fine tune LLM với kiến thức
mới}\label{nghiuxean-cux1ee9u-quy-truxecnh-fine-tune-llm-vux1edbi-kiux1ebfn-thux1ee9c-mux1edbi}

Fine-tuning là quá trình điều chỉnh một mô hình ngôn ngữ lớn (LLM) đã
được huấn luyện trước để học thêm các kiến thức mới hoặc thích nghi với
một tác vụ cụ thể. Quá trình này giúp tiết kiệm thời gian và tài nguyên
so với việc huấn luyện từ đầu, đồng thời tận dụng được khả năng hiểu
ngôn ngữ đã có sẵn của mô hình.

\paragraph{Chuẩn bị dữ
liệu}\label{chuux1ea9n-bux1ecb-dux1eef-liux1ec7u-1}

Dữ liệu fine-tune cần phù hợp với mục tiêu ứng dụng: có thể là hội
thoại, câu hỏi--trả lời, dữ liệu kỹ thuật, hoặc tài liệu chuyên ngành.
Dữ liệu cần được định dạng đúng chuẩn (thường theo dạng
prompt--response) và có chất lượng cao để đảm bảo hiệu quả fine-tune.

\includegraphics[width=4.97396in,height=2.62247in]{media/image12.png}

\emph{\textbf{Hình 3.2.3.1. Dữ liệu dạng hỏi/trả lời để fine tune LLM}}

\emph{Nguồn: A Dataset for LLM Question Answering with External Tools,
Yuchen Zhuang. Trích từ: https://neurips.cc/virtual/2023/poster/73466}

\paragraph{Chuẩn bị hardware}\label{chuux1ea9n-bux1ecb-hardware-1}

Fine-tune yêu cầu phần cứng nhẹ hơn so với huấn luyện từ đầu nhưng vẫn
cần GPU có dung lượng bộ nhớ lớn (tối thiểu 16GB VRAM cho mô hình vài tỷ
tham số). Các kỹ thuật như LoRA (Low-Rank Adaptation) giúp giảm chi phí
phần cứng khi fine-tune.

\includegraphics[width=4.35573in,height=2.45282in]{media/image10.png}

\emph{\textbf{Hình 3.2.3.2. NVIDIA GPU 4090}}

\emph{Nguồn: GeForce RTX 4090, Nvidia. Trích từ:
https://www.nvidia.com/en-gb/geforce/graphics-cards/40-series/rtx-4090/}

\paragraph{Lựa chọn mô hình
LLM}\label{lux1ef1a-chux1ecdn-muxf4-huxecnh-llm}

Chọn mô hình pre-trained phù hợp với kích thước và yêu cầu ứng dụng. Các
mô hình phổ biến để fine-tune hiện nay gồm: LLaMA 3, Mistral, hoặc
Phi-2. Việc lựa chọn mô hình phụ thuộc vào tài nguyên sẵn có, và độ hiệu
quả trên tác vụ mục tiêu.

\includegraphics[width=4.14063in,height=1.91974in]{media/image13.png}

\emph{\textbf{Hình 3.2.3.3. Mô hình Llama của Meta rất phổ biến khi fine
tune LLM.}}

\emph{Nguồn: LLaMA: Meta's Open-Source Rival to Google and OpenAI,
Ankur. Trích từ:
https://www.ankursnewsletter.com/p/llama-metas-open-source-rival-to}

\paragraph{Fine tune LLM}\label{fine-tune-llm}

Thực hiện fine-tune mô hình bằng phương pháp LoRA (Low-Rank Adaptation)
thông qua các thư viện như Hugging Face Transformers, PEFT, hoặc tích
hợp với DeepSpeed để tối ưu hiệu suất. Trong quá trình huấn luyện, theo
dõi các chỉ số như loss và perplexity nhằm đánh giá mức độ học của mô
hình. Có thể áp dụng thêm các kỹ thuật kiểm soát như early stopping và
gradient checkpointing để tiết kiệm tài nguyên và ngăn hiện tượng
overfitting.

\includegraphics[width=4.92576in,height=2.26563in]{media/image14.png}

\emph{\textbf{Hình 3.2.3.4. Sử dụng LoRA khi fine tune LLM.}}

\emph{Nguồn: Low Rank Adaptation (LoRA): Efficient Fine-Tuning for
Real-World AI Applications, Pranjal Khadka. Trích từ:
https://pub.towardsai.net/low-rank-adaptation-lora-fedf37b92026}

\subsection{Nghiên cứu AI Agent}\label{nghiuxean-cux1ee9u-ai-agent}

AI Agent là hệ thống trí tuệ nhân tạo có khả năng tương tác với con
người và môi trường một cách linh hoạt, chủ động, và có mục tiêu rõ
ràng. Trong bối cảnh LLM, AI Agent được thiết kế để mở rộng khả năng của
mô hình ngôn ngữ bằng cách thêm kiến thức mới, tương tác với phần mềm,
và thực hiện chuỗi hành động phức tạp.

\includegraphics[width=3.27877in,height=3.46297in]{media/image15.png}

\emph{\textbf{Hình 3.3. Kiến trúc tổng quan của AI Agent}}

\emph{Nguồn: Agent 101 hơ to build your first ai agent in 30 minutes,
Anmol Baranwal.Trích từ:
https://www.copilotkit.ai/blog/agents-101-how-to-build-your-first-ai-agent-in-30-minutes}

\subsubsection{\texorpdfstring{\textbf{Nghiên cứu cách thêm kiến thức
mới vào cho
LLM}}{Nghiên cứu cách thêm kiến thức mới vào cho LLM}}\label{nghiuxean-cux1ee9u-cuxe1ch-thuxeam-kiux1ebfn-thux1ee9c-mux1edbi-vuxe0o-cho-llm}

\paragraph{Retrieval Augmented Generation
(RAG)}\label{retrieval-augmented-generation-rag}

RAG là kỹ thuật giúp LLM truy cập kiến thức bên ngoài thông qua truy vấn
tìm kiếm tài liệu, sau đó kết hợp nội dung đó để tạo câu trả lời. Thay
vì ``nhồi'' tất cả kiến thức vào mô hình qua fine-tune, RAG cho phép giữ
mô hình gọn nhẹ, dễ cập nhật kiến thức. Các thành phần chính trong RAG
gồm: vector store, retriever, và mô hình LLM.

\includegraphics[width=5.99479in,height=2.89in]{media/image16.png}

\emph{\textbf{Hình 3.3.1.1. Lấy những dữ liệu bên ngoài đưa vào trong
prompt}}

\emph{Nguồn: RAG from scratch. Trích từ:
https://github.com/langchain-ai/rag-from-scratch}

\subsubsection{\texorpdfstring{\textbf{Nghiên cứu cách giúp LLM tương
tác với các phần mềm khác theo thời gian
thực}}{Nghiên cứu cách giúp LLM tương tác với các phần mềm khác theo thời gian thực}}\label{nghiuxean-cux1ee9u-cuxe1ch-giuxfap-llm-tux1b0ux1a1ng-tuxe1c-vux1edbi-cuxe1c-phux1ea7n-mux1ec1m-khuxe1c-theo-thux1eddi-gian-thux1ef1c}

\paragraph{AI use tools}\label{ai-use-tools}

Để giúp LLM không chỉ ``nói'' mà còn ``hành động'', một AI Agent cần khả
năng gọi và sử dụng công cụ (tool usage). Ví dụ: gọi API, truy xuất cơ
sở dữ liệu, gửi email, thực hiện tìm kiếm,... Điều này thường được tích
hợp qua các ``tool executor'' trong hệ thống, cho phép LLM viết và gọi
lệnh Python, HTTP request hoặc các đoạn mã shell để đạt mục tiêu cụ thể.

\includegraphics[width=5.77541in,height=3.46332in]{media/image17.png}

\emph{\textbf{Hình 3.3.2.1. AI Agent sử dụng tool ngoài để đăng bài lên
twitter.}}

\emph{Nguồn: AI agentic workflows: a practical guide for n8n automation,
Yulia Dmitrievna. Trích từ: https://blog.n8n.io/ai-agentic-workflows/}

\subsubsection{\texorpdfstring{\textbf{Khảo sát các AI Agent opensource
hiện
có}}{Khảo sát các AI Agent opensource hiện có}}\label{khux1ea3o-suxe1t-cuxe1c-ai-agent-opensource-hiux1ec7n-cuxf3}

\paragraph{Langchain}\label{langchain}

LangChain là framework mã nguồn mở giúp xây dựng AI Agent sử dụng LLM,
hỗ trợ kết nối dữ liệu, công cụ, bộ nhớ, và thực hiện tác vụ đa bước.
LangChain phổ biến vì cấu trúc rõ ràng, thư viện phong phú và dễ mở
rộng.

\includegraphics[width=3.75145in,height=1.89947in]{media/image18.png}

\emph{\textbf{Hình 3.3.3.1. Langchain logo}}

\emph{Nguồn: Empowering LLMs with LangChain Framework, Debaprasann
Bhroi. Trích từ:
https://blog.gopenai.com/empowering-llms-with-langchain-framework-16c7a0785c2f}

\paragraph{Langgraph}\label{langgraph}

LangGraph là một thư viện mở rộng của LangChain, cho phép xây dựng các
quy trình logic phức tạp theo mô hình đồ thị trạng thái (stateful
graph). Rất phù hợp với các AI Agent cần luồng xử lý rõ ràng, có thể
tuần tự hoặc song song.

\includegraphics[width=3.41302in,height=3.42077in]{media/image19.png}

\emph{\textbf{Hình 2.13. LangGraph logo.}}

\emph{Nguồn: Building graphs for chatbots. Trích từ:
https://campus.datacamp.com/courses/designing-agentic-systems-with-langchain/building-chatbots-with-langgraph?ex=1}

\paragraph{AutoGen}\label{autogen}

AutoGen được phát triển bởi Microsoft, cho phép tạo AI Agent có thể phối
hợp nhiều LLM với nhau theo cách đa tác vụ. Hệ thống này hỗ trợ
``multi-agent collaboration'', nơi các agent đóng vai trò khác nhau
(planner, executor, researcher) để cùng giải quyết một nhiệm vụ.

\includegraphics[width=2.34375in,height=2.34375in]{media/image20.png}

\emph{\textbf{Hình 3.3.3.3. AutoGen logo.}}

\emph{Nguồn: Microsoft's Agentic Frameworks: AutoGen and Semantic
Kernel, Friederike Niedtner. Trích từ:
https://devblogs.microsoft.com/autogen/microsofts-agentic-frameworks-autogen-and-semantic-kernel/}

\paragraph{Semantic Kernel}\label{semantic-kernel}

Semantic Kernel là framework từ Microsoft giúp tích hợp LLM với các công
cụ và tác vụ trong ứng dụng thực tế. Nó hỗ trợ memory (ghi nhớ ngữ
cảnh), planning, và tích hợp công cụ theo mô hình plugin, rất phù hợp để
xây dựng AI Agent doanh nghiệp.

\includegraphics[width=2.34375in,height=2.34375in]{media/image21.png}

\emph{\textbf{Hình 3.3.3.4. Semantic Kernel logo.}}

\emph{Nguồn: Semantic Kernel's new icon and the art of teamwork,
Veronica. Nguồn:
https://devblogs.microsoft.com/semantic-kernel/semantic-kernels-new-icon-and-the-art-of-teamwork/}

\paragraph{CrewAI}\label{crewai}

CrewAI là một thư viện tập trung vào việc xây dựng các AI Agent hoạt
động theo nhóm (crew), mỗi agent có vai trò riêng biệt và phối hợp để
hoàn thành một tác vụ lớn. Mô hình này mô phỏng tổ chức nhóm người làm
việc, với khả năng lập kế hoạch và cộng tác.

\includegraphics[width=3.56927in,height=1.12366in]{media/image22.png}

\emph{\textbf{Hình 3.3.3.5. Crewai logo.}}

\emph{Nguồn: CrewAI Launches Multi-Agentic Platform to Deliver on the
Promise of Generative AI for Enterprise, Daniel D. Gutinerrez. Trích từ:
https://radicaldatascience.wordpress.com/2024/10/22/crewai-launches-multi-agentic-platform-to-deliver-on-the-promise-of-generative-ai-for-enterprise/}

\paragraph{Langflow}\label{langflow}

Langflow là công cụ giao diện trực quan dạng kéo-thả (no-code/low-code)
giúp thiết kế AI Agent dựa trên LangChain. Người dùng có thể xây dựng
pipeline logic với LLM, công cụ, bộ nhớ,... mà không cần viết quá nhiều
mã.

\includegraphics[width=4.25677in,height=1.26904in]{media/image23.png}

\emph{\textbf{Hình 3.3.3.6.1. Langfow logo.}}

\emph{Nguồn: Langflow Reimagined: Streamlined for Even More
Productivity. Trích từ:
https://www.datastax.com/blog/langflow-brand-and-ui-reimagined-with-1-1-release}

\includegraphics[width=6.26772in,height=3.63889in]{media/image24.png}

\emph{\textbf{Hình 3.3.3.6.2. Drag \& drop để tạo workflow trong
langflow.}}

\emph{Nguồn: A Beginner's Guide to Building Agents in Langfloư. Trích từ
: https://www.datastax.com/blog/guide-to-building-agents-in-langflow}

\begin{enumerate}
\def\labelenumi{\arabic{enumi}.}
\setcounter{enumi}{6}
\item
  \textbf{n8n}
\end{enumerate}

n8n là một công cụ workflow automation mã nguồn mở, cho phép xây dựng
quy trình tự động hóa bằng giao diện kéo--thả. Trong bối cảnh AI Agent,
n8n có thể được sử dụng như một ``orchestrator'' -- giúp kết nối LLM với
các API, hệ thống backend, hoặc dịch vụ bên ngoài. n8n hỗ trợ nhiều tích
hợp sẵn như Google Sheets, Slack, Webhook, HTTP Request,\ldots, giúp AI
Agent thực hiện hành động thật ngoài môi trường LLM như gửi email, truy
xuất dữ liệu, hoặc báo cáo kết quả.

\includegraphics[width=3.61419in,height=1.45733in]{media/image25.png}

\emph{\textbf{Hình 3.3.3.7.1. n8n logo.}}

\emph{Nguồn: Khám Phá N8n Workflow và Ứng Dụng Trong Tự Động Hóa AI
2025. Trích từ:
https://kb.pavietnam.vn/n8n-workflow-ung-dung-ai-2025.html}

\includegraphics[width=4.93906in,height=2.53235in]{media/image26.png}

\emph{\textbf{Hình 3.3.3.7.2. Drag \& drop để tạo flow trong n8n.}}

\emph{Nguồn: Building AI Agents with n8n: Low-Code approach to AI
workflows. Trích từ:
https://www.pondhouse-data.com/blog/ai-agents-with-n8n}

\paragraph{RAGFlow}\label{ragflow}

Ragflow là một framework mới nổi tập trung vào việc xây dựng pipeline
RAG (Retrieval-Augmented Generation) dễ dàng hơn, thông qua cấu trúc
mô-đun và có thể cắm thêm công cụ. Ragflow hỗ trợ các thành phần như:
document loaders, vector stores, retriever, chunker và LLM. Nó giúp xây
dựng AI Agent có khả năng truy cập tài liệu bên ngoài, cập nhật kiến
thức liên tục, và tương tác với cơ sở tri thức của doanh nghiệp một cách
trực quan và dễ cấu hình.

\includegraphics[width=2.30885in,height=1.47801in]{media/image27.png}

\emph{\textbf{Hình 3.3.3.8.1. RAGFlow logo.}}

\emph{Nguồn: Meet RAGFlow: An Open-Source RAG (Retrieval-Augmented
Generation) Engine Based on Deep Document Understanding. Trích từ:
https://www.marktechpost.com/2024/04/06/meet-ragflow-an-open-source-rag-retrieval-augmented-generation-engine-based-on-deep-document-understanding/}

\includegraphics[width=4.26198in,height=2.55355in]{media/image28.png}

\emph{\textbf{Hình 3.3.3.8.2. Drag \& drop để tạo flow trong RAGflow.}}

\emph{Nguồn: RAGFlow Enters Agentic Era, Yingfeng Zhang. Trích từ:
\href{https://ragflow.io/blog/ragflow-enters-agentic-era}{\ul{https://ragflow.io/blog/ragflow-enters-agentic-era}}}

\subsection{Agentic system}\label{agentic-system}

\subsubsection{Orchestration}\label{orchestration}

Agentic Orchestration là khái niệm trong lĩnh vực trí tuệ nhân tạo và hệ
thống đa tác tử (multi-agent systems), chỉ việc điều phối và quản lý
nhiều tác nhân (agents) tự chủ để đạt được các mục tiêu phức tạp. Trong
cơ chế này:

\begin{itemize}
\item
  Agents là các thực thể có khả năng ra quyết định độc lập, thực hiện
  hành động dựa trên quan sát và mục tiêu riêng.
\item
  Orchestration là quá trình kiểm soát, phối hợp hoặc tối ưu hóa hành vi
  của các agents để đạt hiệu quả chung, tránh xung đột và khai thác tốt
  khả năng của từng agent.
\item
  Agentic Orchestration nhấn mạnh vào tính tự chủ có điều hướng, tức mỗi
  agent vẫn duy trì quyền ra quyết định riêng nhưng được điều chỉnh theo
  chiến lược tổng thể.
\end{itemize}

\begin{itemize}
\item
  AI Agents trong doanh nghiệp tự động hóa các quy trình phức tạp.
\item
  Các hệ thống đa tác vụ như trợ lý ảo thông minh, nơi nhiều agents phối
  hợp để hoàn thành các nhiệm vụ khác nhau.
\item
  Robot swarm, game AI, hoặc các hệ thống phân tán cần hợp tác và điều
  phối động.
\end{itemize}

\includegraphics[width=6.26772in,height=3.19444in]{media/image29.png}

\emph{\textbf{Hình 3.4.1. Agentic Orchestration.}}

\emph{Nguồn: ``Think Big, Act Small: What Is Agentic AI and AI
Distillation?'', concentrix. Trích từ :
\href{https://www.concentrix.com/insights/blog/what-is-agentic-ai-and-ai-distillation/}{\ul{https://www.concentrix.com/insights/blog/what-is-agentic-ai-and-ai-distillation/}}}

\subsubsection{Planning Agent}\label{planning-agent}

Planning Agent là một loại tác nhân (agent) trong trí tuệ nhân tạo có
khả năng lập kế hoạch hành động trước khi thực hiện, thay vì chỉ phản
ứng tức thời với môi trường. Điểm đặc trưng của Planning Agent là nó sử
dụng thông tin về trạng thái hiện tại, mục tiêu và các hành động khả thi
để tính toán một chuỗi hành động tối ưu nhằm đạt được mục tiêu cuối
cùng.

Các thành phần chính:

\begin{itemize}
\item
  Môi trường (Environment): nơi agent thực hiện hành động và nhận phản
  hồi.
\item
  Trạng thái (State): mô tả thông tin hiện tại của môi trường.
\item
  Hành động (Action): tập các hành động mà agent có thể thực hiện.
\item
  Kế hoạch (Plan): chuỗi hành động được lập trước để đạt mục tiêu.
\item
  Mục tiêu (Goal): trạng thái mà agent muốn đạt tới.
\end{itemize}

Ứng dụng:

\begin{itemize}
\item
  Robot di chuyển hoặc thực hiện nhiệm vụ phức tạp.\\
  AI trong game để đưa ra chiến lược tối ưu.
\item
  Hệ thống trợ lý thông minh thực hiện đa tác vụ (multi-tasking AI
  Agent).
\end{itemize}

\includegraphics[width=6.26772in,height=4.19444in]{media/image30.png}

\emph{\textbf{Hình 2.1.2. Planning agent.}}

\emph{Nguồn: ``Plan-and-Execute Agents'', Langchain. Trích từ :
https://blog.langchain.com/planning-agents/}

\subsubsection{Specialist AI Agent}\label{specialist-ai-agent}

Specialist Agent là một loại tác nhân (agent) trong trí tuệ nhân tạo
được thiết kế tập trung vào một lĩnh vực hoặc nhiệm vụ cụ thể, với khả
năng thực hiện chuyên sâu hơn so với các agent tổng quát. Thay vì cố
gắng xử lý nhiều loại nhiệm vụ khác nhau, Specialist Agent được tối ưu
để hoàn thành hiệu quả một hoặc vài loại công việc chuyên biệt.

Đặc điểm chính:

\begin{itemize}
\item
  Chuyên môn hóa: agent được huấn luyện hoặc lập trình để xử lý một loại
  dữ liệu, nhiệm vụ, hoặc domain cụ thể.
\item
  Hiệu quả cao: nhờ tập trung vào phạm vi hẹp, agent thường đưa ra quyết
  định chính xác và nhanh hơn.
\item
  Tương tác với các agent khác: trong các hệ thống đa-agent, Specialist
  Agents có thể phối hợp với Planning Agents hoặc Generalist Agents để
  thực hiện các nhiệm vụ phức tạp.
\end{itemize}

Ứng dụng:

\begin{itemize}
\item
  AI phân tích văn bản y tế chuyên sâu.
\item
  Chatbot chuyên trả lời về luật pháp hoặc tài chính.
\item
  Hệ thống robot thực hiện một nhiệm vụ cụ thể như hàn hoặc lắp ráp
  trong dây chuyền sản xuất.
\end{itemize}

\includegraphics[width=2.86458in,height=1.90625in]{media/image31.png}

\emph{\textbf{Hình 3.4.3. Specialist Agent.}}

\emph{Nguồn: ``What is a SEO Specialist?'', Analytics Vidhya. Trích từ :
https://relevanceai.com/agent-templates-roles/seo-specialist-ai-agents}
